\documentclass[../main.tex]{subfiles}
\begin{document}

\section{Introduction}
\label{sec:introduction}

Systematic reviews and meta-analyses occupy the apex of the evidence hierarchy in
evidence-based medicine \citep{Higgins2019}.
Their validity depends critically on the quality of the primary studies included:
a meta-analysis that aggregates trials with uncontrolled confounding or selective
outcome reporting inherits---and can amplify---those methodological flaws.
The Cochrane Risk of Bias~1.0 tool (RoB~1.0) \citep{Higgins2011} was developed
precisely to render this quality appraisal explicit and reproducible.
By requiring reviewers to judge each included trial across seven pre-specified
methodological domains, RoB~1.0 transforms a subjective quality impression into a
structured, auditable record.

Despite this formalisation, RoB assessment remains a major bottleneck in systematic
review production.
Each trial must be read by two independent reviewers, who compare judgements and
resolve disagreements through discussion or adjudication.
For a review covering dozens of trials, this process routinely demands tens to
hundreds of person-hours.
Inter-rater reliability, even among trained reviewers, is moderate at best
\citep{Konsgen2020, ArmijoOlivo2014}---a consequence of genuine ambiguity in
reporting, domain-specific expertise requirements, and the epistemic character of
the \emph{Unclear Risk} category, which denotes insufficient information rather
than a graded risk level.
These limitations have motivated sustained interest in automating, or at least
augmenting, the RoB assessment workflow.

Early automation efforts relied on rule-based systems and supervised machine learning
trained on sentence-level text fragments \citep{Marshall2016}.
Such approaches demonstrated that natural language processing could identify
methodological descriptions in trial reports with above-chance accuracy, but they
were constrained by limited contextual understanding, sensitivity to reporting style,
and the need for substantial labelled training data for each domain.

Large language models (LLMs) trained on diverse corpora with instruction-following
capabilities represent a qualitatively different paradigm
\citep{Lieberum2025, Scherbakov2025}.
Without task-specific fine-tuning, general-purpose LLMs can be prompted to read a
full trial report, locate domain-relevant passages, and produce structured judgements
with free-text justifications---closely mimicking the expert reviewer workflow.
A prompt template can address all seven RoB~1.0 domains either sequentially---one
domain per inference call---or simultaneously in a single call; whether this
choice affects assessment quality is an open empirical question.
As frontier LLMs have grown more capable and their inference costs have fallen,
evaluating their effectiveness for RoB assessment has become both feasible and
practically important.
A recent preregistered comparative study applied four contemporary LLMs to
100~RCTs from Cochrane reviews under the RoB~1.0 instrument, reporting
interobserver $\kappa$ of 0.27--0.39 and concluding that none reached the
reliability threshold required for autonomous use \citep{Nyrhi2026}.
Whether this shortfall is uniform across the RoB~1.0 structure---or whether
reliable assistance is already viable for specific domains and intervention types
while systematic failures persist elsewhere---has direct implications for how
LLMs can be integrated into real screening workflows.

The present study addresses this question with a systematic comparison of eight
state-of-the-art LLMs on zero-shot, single-agent RoB~1.0 assessment applied to
227~randomised controlled trials drawn from 123~Cochrane systematic reviews
published in~2024.
All models were evaluated on a common benchmark under identical conditions: the same
prompt, the same fully text-based study representations produced by a standardised
PDF-to-Markdown pipeline, and the same ground-truth human annotations.
Beyond aggregate accuracy and Macro-F1, we report Cohen's~$\kappa$, Gwet's~AC1,
and per-class F1---placing $\kappa$ and AC1 side by side to expose reliability
estimates that diverge under class imbalance.
Domain- and intervention-stratified analyses identify where systematic biases arise
and locate the improvement opportunities most likely to yield practical gains for
assisted screening.
Additionally, we compare two inference strategies---per-domain and
single-call---to assess whether querying all seven domains simultaneously
affects performance relative to sequential per-domain calls.

\end{document}
